% GENERATED FILE. Edit canonical lessons, then run npm run generate:books.
% canonical-source-hash: fnv1a64:af4642a8318acd88
% canonical-lessons: LA-C502-iocus, LA-C502-iocor, LA-C502-iocosus, LA-C502-ridiculus, LA-C502-facetus

\chapter{Jokes and wit}
\label{ch:la-jokes-and-wit}
\hlchaptermodality{502}

\begin{chapteropening}
\textbf{By the end of this chapter:} \emph{I can talk about jokes and witty people.}

Complete the last lesson of chapter 502: I can talk about jokes and witty people.
\end{chapteropening}

\section[iocus, iocī]{iocus, iocī — a joke, jest}

\label{lesson:LA-C502-iocus}



\begin{quote}
Before the new one: say the Latin for calm, then the Latin for free from care.
\end{quote}

\subsection*{You'll want to know: iocus, iocī}
\textbf{iocus, iocī} — \textquotedblleft{}a joke, jest\textquotedblright{}.

\begin{quote}
Pater iocum nārrat, et omnēs rīdent.
\end{quote}

\begin{quote}
\emph{Father tells a joke, and everyone laughs.}
\end{quote}

Masculine, second declension; the plural is often \emph{ioca}. \emph{Per iocum} — as a joke. English \emph{joke} and \emph{jocular}.

\begin{cousinweb}
\textbf{iocus} is inherited through Proto-Italic from a Proto-Indo-European word that EtymDB traces to a root meaning \textquotedblleft{}to speak\textquotedblright{}, with Welsh \emph{iaith} \textquotedblleft{}language\textquotedblright{} as a cousin; Lewis \& Short is less sure. English \emph{joke} and \emph{jocular} come from it.
\end{cousinweb}

\begin{grammarlens}[title={What you've built}]
One more word for telling, explaining and describing at B1.
\end{grammarlens}

\subsection*{Your turn}
\begin{itemize}
\raggedright
  \item \emph{Say it:} \emph{iocus, iocī}
  \item \emph{Say it:} \emph{iocus, iocī}, once more
  \item \emph{Say it:} a sentence of your own with \emph{iocus, iocī}
  \item \emph{Recall:} say the Latin for calm, then the Latin for free from care, then say \emph{iocus, iocī} again
\end{itemize}

\begin{tcolorbox}[breakable,colback=teal!4,colframe=teal!35!black,title={Before you move on}]
What does iocus, iocī mean? (\textquotedblleft{}A joke, jest\textquotedblright{}.) Say it once more.
\end{tcolorbox}
\section[iocor, iocārī]{iocor, iocārī — to joke, jest}

\label{lesson:LA-C502-iocor}



\begin{quote}
Before the new one: say the Latin for free from care, then the Latin for a joke.
\end{quote}

\subsection*{You'll want to know: iocor, iocārī}
\textbf{iocor, iocārī} — \textquotedblleft{}to joke, jest\textquotedblright{}.

\begin{quote}
Servī in culīnā iocantur.
\end{quote}

\begin{quote}
\emph{The slaves joke in the kitchen.}
\end{quote}

Deponent: \emph{iocor, iocārī, iocātus sum}.

\begin{cousinweb}
\textbf{iocor} is built on \emph{iocus} \textquotedblleft{}a joke\textquotedblright{} (Lewis \& Short; EtymDB).
\end{cousinweb}

\begin{grammarlens}[title={What you've built}]
One more word for telling, explaining and describing at B1.
\end{grammarlens}

\subsection*{Your turn}
\begin{itemize}
\raggedright
  \item \emph{Say it:} \emph{iocor, iocārī}
  \item \emph{Say it:} \emph{iocor, iocārī}, once more
  \item \emph{Say it:} a sentence of your own with \emph{iocor, iocārī}
  \item \emph{Recall:} say the Latin for free from care, then the Latin for a joke, then say \emph{iocor, iocārī} again
\end{itemize}

\begin{tcolorbox}[breakable,colback=teal!4,colframe=teal!35!black,title={Before you move on}]
What does iocor, iocārī mean? (\textquotedblleft{}To joke, jest\textquotedblright{}.) Say it once more.
\end{tcolorbox}
\section[iocōsus, iocōsa, iocōsum]{iocōsus, iocōsa, iocōsum — humorous, playful}

\label{lesson:LA-C502-iocosus}



\begin{quote}
Before the new one: say the Latin for a joke, then the Latin for to joke.
\end{quote}

\subsection*{You'll want to know: iocōsus, iocōsa, iocōsum}
\textbf{iocōsus, iocōsa, iocōsum} — \textquotedblleft{}humorous, playful\textquotedblright{}.

\begin{quote}
Avus meus senex iocōsus est.
\end{quote}

\begin{quote}
\emph{My grandfather is a humorous old man.}
\end{quote}

Built on \emph{iocus}. English \emph{jocose}.

\begin{cousinweb}
\textbf{iocōsus} is built on \emph{iocus} \textquotedblleft{}a joke\textquotedblright{} with \emph{-ōsus} \textquotedblleft{}full of\textquotedblright{} (Lewis \& Short; EtymDB). English \emph{jocose} comes from it.
\end{cousinweb}

\begin{grammarlens}[title={What you've built}]
One more word for telling, explaining and describing at B1.
\end{grammarlens}

\subsection*{Your turn}
\begin{itemize}
\raggedright
  \item \emph{Say it:} \emph{iocōsus, iocōsa, iocōsum}
  \item \emph{Say it:} \emph{iocōsus, iocōsa, iocōsum}, once more
  \item \emph{Say it:} a sentence of your own with \emph{iocōsus, iocōsa, iocōsum}
  \item \emph{Recall:} say the Latin for a joke, then the Latin for to joke, then say \emph{iocōsus, iocōsa, iocōsum} again
\end{itemize}

\begin{tcolorbox}[breakable,colback=teal!4,colframe=teal!35!black,title={Before you move on}]
What does iocōsus, iocōsa, iocōsum mean? (\textquotedblleft{}Humorous, playful\textquotedblright{}.) Say it once more.
\end{tcolorbox}
\section[rīdiculus, rīdicula, rīdiculum]{rīdiculus, rīdicula, rīdiculum — funny, laughable; absurd}

\label{lesson:LA-C502-ridiculus}



\begin{quote}
Before the new one: say the Latin for to joke, then the Latin for humorous.
\end{quote}

\subsection*{You'll want to know: rīdiculus, rīdicula, rīdiculum}
\textbf{rīdiculus, rīdicula, rīdiculum} — \textquotedblleft{}funny, laughable; absurd\textquotedblright{}.

\begin{quote}
Histriō in theātrō rīdiculus est.
\end{quote}

\begin{quote}
\emph{The actor is funny in the theatre.}
\end{quote}

From \emph{rīdeō} — to laugh. Funny in a good sense, absurd in a bad one. English \emph{ridiculous}.

\begin{cousinweb}
\textbf{rīdiculus} is built on \emph{rīdeō} \textquotedblleft{}to laugh\textquotedblright{} (Lewis \& Short; EtymDB). English \emph{ridiculous} and \emph{ridicule} come from it.
\end{cousinweb}

\begin{grammarlens}[title={What you've built}]
One more word for telling, explaining and describing at B1.
\end{grammarlens}

\subsection*{Your turn}
\begin{itemize}
\raggedright
  \item \emph{Say it:} \emph{rīdiculus, rīdicula, rīdiculum}
  \item \emph{Say it:} \emph{rīdiculus, rīdicula, rīdiculum}, once more
  \item \emph{Say it:} a sentence of your own with \emph{rīdiculus, rīdicula, rīdiculum}
  \item \emph{Recall:} say the Latin for to joke, then the Latin for humorous, then say \emph{rīdiculus, rīdicula, rīdiculum} again
\end{itemize}

\begin{tcolorbox}[breakable,colback=teal!4,colframe=teal!35!black,title={Before you move on}]
What does rīdiculus, rīdicula, rīdiculum mean? (\textquotedblleft{}Funny, laughable; absurd\textquotedblright{}.) Say it once more.
\end{tcolorbox}
\section[facētus, facēta, facētum]{facētus, facēta, facētum — witty, clever}

\label{lesson:LA-C502-facetus}



\begin{quote}
Before the new one: say the Latin for humorous, then the Latin for funny.
\end{quote}

\subsection*{You'll want to know: facētus, facēta, facētum}
\textbf{facētus, facēta, facētum} — \textquotedblleft{}witty, clever\textquotedblright{}.

\begin{quote}
Hospes facētus multa rīdicula dīcit.
\end{quote}

\begin{quote}
\emph{The witty guest says many funny things.}
\end{quote}

Elegant, good-humoured wit. English \emph{facetious} comes from \emph{facētiae} — jokes.

\begin{cousinweb}
\textbf{facētus} is a native Latin adjective; Lewis \& Short and EtymDB link it with \emph{fax} \textquotedblleft{}torch\textquotedblright{} and a root meaning \textquotedblleft{}to shine\textquotedblright{}, but the origin is not settled. English \emph{facetious} comes from \emph{facētiae} \textquotedblleft{}jokes\textquotedblright{}, built on it.
\end{cousinweb}

\begin{grammarlens}[title={What you've built}]
One more word for telling, explaining and describing at B1.
\end{grammarlens}

\subsection*{Your turn}
\begin{itemize}
\raggedright
  \item \emph{Say it:} \emph{facētus, facēta, facētum}
  \item \emph{Say it:} \emph{facētus, facēta, facētum}, once more
  \item \emph{Say it:} a sentence of your own with \emph{facētus, facēta, facētum}
  \item \emph{Recall:} say the Latin for humorous, then the Latin for funny, then say \emph{facētus, facēta, facētum} again
\end{itemize}

\begin{tcolorbox}[breakable,colback=teal!4,colframe=teal!35!black,title={Before you move on}]
What does facētus, facēta, facētum mean? (\textquotedblleft{}Witty, clever\textquotedblright{}.) Say it once more.
\end{tcolorbox}
